\section{Conclusions}
\label{sec:conclusions}

\begin{enumerate}[leftmargin=1.6em,itemsep=4pt,topsep=4pt,parsep=0pt]

\item \emph{What is new.} In the millimetre band the record for nearby stars
is two measurements of one spectral line: \citet{Steffes1994} and
\citet{Mauersberger1996} each covered no more than \SbPriorBandMHz\,MHz at
narrowband resolution around the \SbPriorFreqGHz\,GHz positronium line, with
no search over drift rate and
no epoch-to-epoch test of persistence. This survey covers at
least \SbPriorRatio{} times that bandwidth, searches a grid of drift rates,
and asks of every event whether it returns (\S\ref{sec:background}).

\item \emph{What was searched.} The carrier search is \NWinA{} fine-channel
windows toward \NSysClassA{} systems within 40\,pc, covering \CvUnionA\,GHz
of unique frequency in \CovNIslands{} disjoint intervals. In \EdrNWinA{} of
those windows the track is long enough for the searched drift to cross a
channel, covering \EdrUnionGHz\,GHz toward \EdrNSys{} systems. A coarse-channel
search for a steady unresolved excess adds \NWinB{} windows toward
\PrSysClassB{} systems. Both come from one pass through \NEB{} public ALMA
execution blocks and together span \SurvFreqLo--\SurvFreqHi\,GHz
(\S\ref{sec:sample}).

\item \emph{How deep.} The median system's best window would have produced a
threshold crossing from a carrier of $\EirpNinetySysMedA$\,W nine times in
ten, and \NSysEfifteen{} systems reach $10^{15}$\,W on that measure. That is
the completeness of a single observation and not of a confirmation
(\S\ref{sec:confirmability}); for \SmcNEstSys{} of the \NSysClassA{} systems
it is carried across from the windows that were injected into and not
measured in their own; and it leaves out atmospheric decorrelation, which a
real carrier suffers and an injected one does not, so the net bias on any one
limit runs either way and Table~\ref{tab:persys} gives it system by system
(\S\ref{sec:pninetysel}).

\item \emph{What was found.} The search produced \EvNCrossRaw{} threshold
crossings: \CkNCrossCenA{} Class~A in the primary census, \PrFlagXrossWord{}
Class~B, and \CeNTailCrossWord{} Class~A in archival-tail blocks. All
\PrFlagXrossWord{} Class~B crossings come from the single execution block
that fails the data-quality criterion of Appendix~\ref{sec:blockquality}, so
the quality-passing data hold \EvNCrossQp{} crossings, every one of them from
the carrier search. \EvNAttrQp{} of those are attributed to a
catalogued molecular transition within $\pm\MaskHalfKms$\,km\,s$^{-1}$ of the
star's velocity, circumstellar or interstellar, and \EvNUnattrQp{} are not.
An attribution is not an identification, and \S\ref{sec:mmdifferent} says which of them carry
evidence beyond it. Of the \EvNUnattrQp{}, \CeNTailCrossWord{} lie in
archival-tail blocks that enter no expectation and \CeTlObsCutWord{} is
atmospheric, leaving \CeObs{} to set against the \CeExp{} these fields' own
control positions predict: a count that large or larger arises by chance with
probability \CeProbNull, or \FrNarrowObs{} against \FrNarrowExp{} at
$P=\FrNarrowProb$ over a narrower field, which drops integrations and so is
not the statistic \S\ref{sec:statistic} defines. Either way this is
consistency and not a demonstration that the crossings are noise, since a
further \CeAddLim{} would have been consistent too, \FrNarrowAddLim{} on the
narrower reading (\S\ref{sec:technosearch}, \S\ref{sec:trigunif}).

\item \emph{What the repeat observations exclude.} Nothing came back. A
transmitter on an orbiting body arrives at a frequency its own motion has
moved, so every unattributed crossing another block covers was searched again
twice over: across that whole covering window at every drift of the grid,
where \WdNExc{} of the \WdNCov{} windows hold a statistic above the trigger
against the \WdExpect{} chance predicts in trials of that size, and at the
single cell the crossing's own discovery predicts, where the largest
statistic anywhere is \RcTMatched{} at the drift fitted at discovery and
\RcTMax{} at any drift of the grid, against a trigger of \StageNullTrig. The
same reading recovers \MkAttrRecur{} of the \MkAttrTested{} line-attributed
crossings it was tested on, each in a different observing session more than a
day away (Table~\ref{tab:confirm}), so the machinery does return a repeat
where one exists. That exclusion is narrow, and the
narrowness belongs to the criterion and not to the archive: a carrier radiating in both
observations whose stellar-frame frequency comes back to within half a native
channel. A transmitter accelerating at the \ZtALosLo{} to
\ZtALosHi\,m\,s$^{-2}$ these crossings themselves show re-enters that
interval only if its motion is periodic and the second epoch catches it at
the matching phase, so this is not a test of a persistent transmitter of any
kind. Where two crossings of one star might have been a single carrier,
\TjNPair{} pairs carry enough measurement to be asked: the searched grid
separates \TjNExcl, and the other \TjNOpen{} are \emph{open} -- neither a
confirmation nor an exclusion, with only the agreement of their number with
chance bearing on them (\S\ref{sec:etacrv}).

\item \emph{What could have been confirmed at all.} A search limit and a
confirmation limit are different statements, and here the second is the
binding one. Of the \NSysClassA{} systems the carrier search covered,
\ZcConfSysDay{} were observed again in another observing session more than a
day later at a frequency the first block covered, and \ZcConfSysYr{} more
than a year later; for the remaining \ZcNoRepSys{} no repeat exists, so no
transmitted power would have made a confirmation possible
(\S\ref{sec:confirmability}). Coverage limits this search, and no depth
of observation could have removed that limit.

\item \emph{Stacking the epochs.} Combining each star's observations in its
own rest frame deepens the flux-density limit by a median factor \SyStkGainMed. The
deepest result is Proxima Centauri's: $\SzProxEirpCensus$\,W from the census
blocks and $\StkProxEirpCal$\,W once the rest of the archive is added, in
coarse \ChanBModeMHz\,MHz channels, so it bounds a steady unresolved excess
and not a drifting carrier (\S\ref{sec:stack}). Stacking buys that depth
with a tolerance of one channel, and \PxPlanet's own orbit carries a
transmitter across \PxNChanWalk{} of them, so the stacked figure applies to
the one body in that system only if the orbit is within
\PxIncMaxDeg$^\circ$ of face-on. The single-epoch limit,
$\PxEirpSingle$\,W, applies at any inclination.

\item \emph{What this does not constrain, and what would change the
answer.} These limits hold only inside the domain actually searched. The
drift grid stops at \DriftCeilLo\,nHz, \AccelSheikhPct{} per cent of the
ceiling recommended for a general carrier search \citep{Sheikh2019}, so an
unbound transmitter, one on a rapidly rotating planet, or one whose orbital
and rotational accelerations together exceed that is not constrained here at
all (\S\ref{sec:excludes}). The obstacle transfers beyond
this archive: channelisation chosen by other observers, and not collecting area,
decides which search is possible -- and three choices a proposal can simply
make would lift it, namely fine channels everywhere, frequencies and targets
chosen against the line forest and not inherited from programmes that
sought it, and repeat visits at one tuning days and years apart
(\S\ref{sec:future}).

\end{enumerate}
